\section{Correspondence to other theories}
\label{sec:correspondence_to_other_theories}

\subsection{Classical compositional-flow limit and interpretation of the multipliers}
\label{sec:classical_compositional_limit}

The preceding equations contain the usual compositional-reservoir structure as a
special limit.  The phase/component source is
%
\begin{equation}
	\dot c_\xi^\alpha=\sum_m\nu_{\xi (m)}^\alpha\dot r_{(m)},
	\label{eq:classical_source_form}
\end{equation}
%
which is the phase-resolved analogue of the standard reactive-transport source.
Restrict attention to the fluid pore phases and write their volume fractions as
%
\begin{equation}
	\phi_f=\phi S_f,
	\qquad
	\sum_{f\in\mc F}S_f=1,
	\label{eq:classical_saturation_identification}
\end{equation}
%
where $\phi$ is porosity and $S_f$ is the saturation of fluid phase $f$.
Substitution in \eqref{eq:component_spatial_mass} gives the mass form of the
phase/component compositional balance,
%
\begin{equation}
	\frac{\partial}{\partial t}\left(\phi S_f\bar\rho_f\eta_f^\alpha\right)
	+\nabla_{\mbf x}\cdot
	\left(\phi S_f\bar\rho_f\eta_f^\alpha\mbf v_f\right)
	=\sum_m\nu_{f (m)}^\alpha\dot r_{(m)} .
	\label{eq:classical_phase_mass_balance}
\end{equation}
%
This is the usual phase/component compositional balance, written per unit mass
rather than per mole.  With mechanical dispersion and molecular diffusion
omitted, it is the mass counterpart of Balhoff's standard compositional balance
\cite[eq.~(7.15b)]{balhoff2022multiphase}.  If the equation is written relative
to the solid skeleton, the phase Darcy flux is
%
\begin{equation}
	\mbf q_f=\phi S_f(\mbf v_f-\mbf v_s).
	\label{eq:classical_darcy_flux_definition}
\end{equation}
%
For the fixed-skeleton case $\mbf v_s=\mbf 0$, so
\eqref{eq:classical_phase_mass_balance} becomes
%
\begin{equation}
	\frac{\partial}{\partial t}\left(\phi S_f\bar\rho_f\eta_f^\alpha\right)
	+\nabla_{\mbf x}\cdot
	\left(\bar\rho_f\eta_f^\alpha\mbf q_f\right)
	=\sum_m\nu_{f (m)}^\alpha\dot r_{(m)} .
	\label{eq:classical_phase_mass_darcy_balance}
\end{equation}
%
The standard Darcy closure follows from
\eqref{eq:standard_darcy_mass_flux_limit} by taking $\mbf v_s=\mbf 0$ and
$\bs\kappa_f=\kappa_{fr}(\{S_g\}_{g\in\mc F})\bs\kappa$:
%
\begin{equation}
	\mbf q_f
	=-\frac{\kappa_{fr}(\{S_g\}_{g\in\mc F})}{\mu_f}
	\bs\kappa
	\left(\nabla_{\mbf x}\bar p_f-\bar\rho_f\mbf g\right),
	\label{eq:classical_darcy_closure}
\end{equation}
%
which is Balhoff's Darcy law \cite[eq.~(178)]{balhoff2022multiphase} in the
present notation.  In the present theory this is not postulated at the
balance-law level: it is a closure that replaces the fluid-phase momentum
equation \eqref{eq:el_mom_f_simplified} after inertia, body interactions, and
conversion-momentum terms have been specialized.

Summing \eqref{eq:classical_phase_mass_darcy_balance} over fluid phases yields
%
\begin{equation}
	\frac{\partial}{\partial t}
	\left(\phi\sum_{f\in\mc F}S_f\bar\rho_f\eta_f^\alpha\right)
	+\nabla_{\mbf x}\cdot
	\left(\sum_{f\in\mc F}\bar\rho_f\eta_f^\alpha\mbf q_f\right)
	=\sum_{f\in\mc F}\sum_m\nu_{f (m)}^\alpha\dot r_{(m)} .
	\label{eq:classical_total_mass_component_balance}
\end{equation}
%
Thus the current balance law already has the standard compositional-flow form
once the reservoir variables $\phi$ and $S_f$ are introduced.  The difference is
that the present stoichiometric coefficients retain the phase label.  A pure
phase transformation is the special case in which the component label is
unchanged and the two nonzero entries satisfy
$\nu_{\xi (m)}^\alpha=-1$ and $\nu_{\zeta (m)}^\alpha=+1$; after summing over
phases, its contribution to the total component source vanishes, but it still
changes the phase saturations and phase compositions.

The remaining algebraic structure is also recovered.  The within-phase
composition constraints are
%
\begin{equation}
	\sum_\alpha\eta_\xi^\alpha=1 ,
	\label{eq:classical_composition_constraints}
\end{equation}
%
The thermodynamic distinction appears in how the phase/component potentials are
obtained.  A traditional compositional simulator supplies an equation of state,
phase-equilibrium rules, capillary-pressure relations, component partitioning,
and stress laws as separate closure ingredients.  In the present formulation the
reversible thermodynamic part is instead generated from the Helmholtz
free-energy structure:
%
\begin{equation}
	\psi_\xi
	\quad\Longrightarrow\quad
	\bar p_\xi,\quad
	\mathfrak{s}_\xi,\quad
	\bs\sigma_s',\quad
	\rho_\xi\left.\p{\psi_\xi}{\eta_\xi^\alpha}\right\vert_{\cdots},\quad
	\frac{\tau_\xi^\alpha}{\eta_\xi^\alpha}.
	\label{eq:helmholtz_generates_reversible_closures}
\end{equation}
%
Thus pressure, entropy, stress, composition derivatives, and normalized
component storage potentials are not independent reversible closures; they are
conjugate derivatives or multipliers associated with the constrained Helmholtz
structure.  Transport and reaction kinetics still require dissipative closures,
such as mobilities, permeabilities, diffusivities, and reaction-rate laws, but
the thermodynamic forces driving those laws are generated by the same
free-energy potential.

Classical flash calculations impose phase-equilibrium closures, commonly written
as fugacity equalities.  In the present notation the reversible transfer of
component $\alpha$ from phase $\xi$ to phase $\zeta$ gives
%
\begin{equation}
	\mu_\zeta^\alpha-\mu_\xi^\alpha=0,
	\label{eq:classical_phase_equilibrium_mu}
\end{equation}
%
which is the chemical-potential form of the fugacity-equality condition used in
isothermal phase equilibrium.  In classical compositional models the
$\mu_\xi^\alpha$ are therefore not solved from an evolution law like
%
\begin{equation}
	\frac{D_\xi\mu_\xi^\alpha}{Dt}
	=
	\frac{1}{2}\mbf v_\xi\cdot\mbf v_\xi
	-\mu_\xi^\alpha
	-\frac{\tau_\xi^\alpha}{\phi_\xi\bar\rho_\xi\eta_\xi^\alpha}.
	\label{eq:mu_interpretation_entropy_closure_unified}
\end{equation}
%
Instead, they must be defined constitutively from the chosen equation of state,
composition variables, and temperature.  The variational formulation identifies
\eqref{eq:mu_interpretation_entropy_closure_unified} as the conversion-potential
evolution law associated with the phase-attached source constraint.

The affinity is the stoichiometric projection of these phase/component
potentials:
%
\begin{equation}
	\mc A_{(m)}
	=-\sum_\xi\sum_\alpha \mu_\xi^\alpha\nu_{\xi (m)}^\alpha .
	\label{eq:classical_affinity_projection}
\end{equation}
%
For an equilibrated mechanism, imposing $\mc A_{(m)}=0$ gives
%
\begin{equation}
	\sum_\xi\sum_\alpha \mu_\xi^\alpha\nu_{\xi (m)}^\alpha=0,
	\label{eq:classical_equilibrium_affinity}
\end{equation}
%
which is the classical statement that the stoichiometric sum of chemical
potentials vanishes.  For a kinetically limited mechanism, such as a slow
fluid--solid reaction, \eqref{eq:classical_affinity_projection} instead supplies
the driving force in a closure
%
\begin{equation}
	\dot r_{(m)}=\mc K_{(m)}(\mc A_{(m)},\mbf\Lambda),
	\qquad \mc A_{(m)}\dot r_{(m)}\ge 0
	\label{eq:classical_kinetic_affinity_closure}
\end{equation}
%
with the sign convention used here.  Thus equilibrium chemistry, kinetic
reactive transport, and phase transformation all use the same stoichiometric
projection; only the closure for $\dot r_{(m)}$ changes.

The momentum equations contain an additional conversion term absent from purely
kinematic Darcy-scale compositional balances,
%
\begin{equation}
	\sum_\alpha\sum_\xi\dot c_\xi^\alpha
	\left(\nabla_{\mbf x}\mu_\xi^\alpha-\mbf v_\xi\right).
	\label{eq:conversion_momentum_interpretation}
\end{equation}
%
The term combines the inertial bookkeeping required when mass is removed from or
added to a moving phase with the gradient force generated by variation of the
current-rate constraint $\dot c_\xi^\alpha=\dot{\mc{C}}_\xi^\alpha/J_\xi$.  Total
mass conservation constrains the component source rates, but the explicit form
above avoids introducing an additional conversion-velocity notation.  These terms
vanish or are hidden in classical reservoir models when phase momenta are not
solved explicitly and the phase velocities are supplied by Darcy laws.

Finally, the remaining multipliers have standard constraint interpretations:
$\lambda$ enforces saturation of the pore volume by the phases,
$\pi_\xi$ enforces $\sum_\alpha\eta_\xi^\alpha=1$, and $\tau_\xi^\alpha$ enforces
the phase/component storage relation.  The novelty of the present formulation is
not the stoichiometric source form itself, but the fact that it is embedded in a
phase-attached variational mechanics in which the Jacobian $J_\xi$ is varied.
That is what exposes the conversion-induced momentum production terms and gives
a direct bridge between mixture mechanics and classical compositional reactive
transport.
